\section{Completions}

\subsection{Exercise - Complete}\label{10.1}

\begin{proof}
Let $A$ be the $\Z$-module $\bigoplus_{i=1}^{\infty} (\Z/p\Z)$. Then, $pA = 0$ and thus $p^nA = 0$ for all $n \ge 1$. Then, we have
$\varprojlim A/p^nA = \varprojlim A = A$.

Let $B$ be the $\Z$-module $\bigoplus_{i=1}^{\infty}(\Z/p^i\Z)$. If $\alpha_i : \Z/p\Z \to \Z/p^i\Z$ is the group homomorphism sending
$1 \mapsto p^{i-1}$, then let $\alpha : A \to B$ be the direct sum of the $\alpha_i$. This is a $\Z$-module homomorphism simply because
it is a homomorphism of abelian groups. Note that $B$ also has a natural $p$-filtration. For $n > 0$, 
$p^n B = \bigoplus_{i=1}^{\infty} p^n \Z / p^i \Z$.
Thus, $\alpha_i^{-1}(p^n B) = \Z/p\Z$ if $i > n$, but $\alpha_i^{-1}(p^n B) = 0$ if $i \le n$. Thus, the $p$-filtration of $A$ induced by
the $p$-filtration of $B$ is $A_n = \bigoplus_{i=1}^n 0 \oplus \bigoplus_{i=n+1}^{\infty} \Z/p\Z$. Hence, 
$A/A_n = \bigoplus_{i=1}^n \Z/p\Z = \prod_{i=1}^n \Z/p\Z$.

Note that the canonical homomorphism $\theta_{n+1} : A/A_{n+1} \to A/A_n$ is the one that projects the first $n$ factors isomorphically onto
$A/A_n = \prod_{i=1}^n \Z/p\Z$, while sending the final $\Z/p\Z$ factor to $0$. As in the proof of Proposition 10.2, we can let 
$M = \prod_{r=1}^{\infty}A/M_r$ and $d : M \to M$ be $(d(a))_n = a_n - \theta_{n+1}(a_{n+1})$. Then, $\varprojlim A/A_n \cong \ker d$. Thus,
$a = (a_n)_n \in \ker d$ if and only if $\theta_{n+1}(a_{n+1}) = a_n$ for all $n$. By the characterization of $\theta_{n+1}$, this means that
the $n$ components of $a_n$ match the first $n$ components of $a_{n+1}$. Since this is true for all $n$, it's clear that 
$\ker d \cong \prod_{i=1}^{\infty}\Z/p\Z$. Hence, $\varprojlim A/A_n \cong \prod_{i=1}^{\infty}\Z/p\Z$.

Letting $C = B/\alpha(A)$, we have an exact sequence
$$A \xlongrightarrow{\alpha} B \to C \to 0.$$
For the $p$-adic completion functor to be right exact, we'd need $\im(\varprojlim A/p^n \to \hat B) = \ker (\hat B \to \hat C)$.
Since we showed $A = \varprojlim A/p^n$, we would need $\im(A \to \hat B) = \ker(\hat B \to \hat C)$.

As defined above, $A_n = \alpha^{-1}(p^nB)$, so we have an exact sequence of $p$-filtrations
$$0 \to A/A_n \to B/p^n B \to C/p^n C \to 0.$$
Since $A/A_{n+1} \to A/A_n$ is clearly surjective for all $n$, we have by Proposition 10.2 that
$$0 \to \varprojlim A/A_n \to \hat B \to \hat C \to 0$$
is exact. Then, $\im(\varprojlim A/A_n \to \hat B) = \ker(\hat B \to \hat C)$. From above, $\varprojlim A/A_n \cong \prod_{i=1}^{\infty}\Z/p\Z$.
Since $A = \bigoplus_{i=1}^{\infty}\Z/p\Z$ is a strict subgroup of $\prod_{i=1}^{\infty}\Z/p\Z$ and $\varprojlim A/A_n \to \hat B$ is injective,
$\im(A \to \hat B)$ is a strict subgroup of $\im(\varprojlim A/A_n \to \hat B) = \ker(\hat B \to \hat C)$, so the sequence
$$\varprojlim A/p^n A \xlongrightarrow{\hat \alpha} \varprojlim B/p^n B \to \varprojlim C/p^n C \to 0$$
is not exact. Hence, $p$-adic completion is not a right exact functor on the category of $\Z$-modules.
\end{proof}

\subsection{Exercise - Complete}\label{10.2}

\begin{proof}
As in Exercise \ref{10.1}, define $A_n := \alpha^{-1}(p^n B)$. We have the short exact sequence
$$0 \to A_n \to A \to A/A_n \to 0.$$
From the proof of Exercise \ref{10.1}, we have $A_n = \bigoplus_{i=n+1}^{\infty} \Z/p\Z$ 
and $A/A_n = \prod_{i=1}^n \Z/p\Z$. We have the natural embeddings $A_{n+1} \to A_n$ 
and natural homomorphisms $A/A_{n+1} \to A/A_n$ that map the $n+1$ component to $0$. Thus, we have commutative 
diagrams with exact rows:
\begin{center}
\begin{tikzcd}
0 \arrow[r] & A_{n+1} \arrow[r]\arrow[d] & A \arrow[r]\arrow[d] & A/A_{n+1} \arrow[r]\arrow[d] & 0 \\
0 \arrow[r] & A_n \arrow[r] & A \arrow[r] & A/A_n \arrow[r] & 0.
\end{tikzcd}
\end{center}
So, 
$$0 \to \{A_n\} \to \{A\} \to \{A/A_n\} \to 0$$
is an exact sequence of inverse systems. By Proposition 10.2, we have an exact sequence
$$0 \to \varprojlim A_n \to \varprojlim A \to \varprojlim A/A_n.$$
Then, $\varprojlim A = A = \bigoplus_{i=1}^{\infty}\Z/p\Z$ and, by Exercise \ref{10.1}, $\varprojlim A/A_n \cong \prod_{i=1}^{\infty}\Z/p\Z$.
Now, clearly $\bigoplus_{i=1}^{\infty}\Z/p\Z \to \prod_{i=1}^{\infty}\Z/p\Z$ is not surjective, so $\varprojlim$ is not right exact.

Let $M = \prod_{k=0}^{\infty} A_k$ ($A_0 = A$) and define $d : M \to M$ via $(d(a))_n = a_n - \theta_{n+1}(a_{n+1})$. Then, we have 
$\varprojlim^1 A_n = \coker d$, by definition. Let $x \in \im(d)$. Then, $x_n = (d(a))_n = a_n - \theta_{n+1}(a_{n+1})$. Since $\theta_{n+1}$ 
is just the identification of $A_{n+1}$ as a subgroup of $A_n$, we can write $\theta_{n+1}(a_{n+1})$ as just $a_{n+1}$ so that $x_n = a_n - a_{n+1}$. 
Thus, $$\sum_{i=0}^n x_i = \sum_{i=0}^n (a_i - a_{i+1}) = a_0 - a_{n+1}$$
for all $n$, so $a_{n+1} = a_0 - \sum_{i=0}^n x_i \in A_{n+1}$ for all $n$. Hence, for any $j \le n+1$, 
$$0 = (a_{n+1})_j = (a_0)_j - \left(\sum_{i=0}^n x_i\right)_j.$$
Since $x_i \in A_i$, $(x_i)_j = 0$ whenever $i \ge j$, so we have $(a_0)_j = \left(\sum_{i=0}^j x_i\right)_j$ for all $j$. 
Now, since $a_0 \in A_0 = \bigoplus_{i=1}^{\infty} \Z/p\Z$, $(a_0)_j = 0$ for sufficiently large $j$. Thus, for sufficiently large $j$,
$$\left(\sum_{i=0}^j x_i\right)_j = (a_0)_j = 0$$
so that $\sum_{i=0}^j (x_i)_j = 0$. Since $x_n \in A_n$ for all $n$, the sum $X = \sum_{i=0}^{\infty} x_i$ is well-defined since it becomes 
stationary in every coordinate. Then, $X_j = \left(\sum_{i=0}^{\infty} x_i\right)_j = \sum_{i=0}^j (x_i)_j$. Thus, for sufficiently large $j$, 
$X_j = 0$. Hence, $X \in A$.

Conversely, suppose for $(x_i)_i \in M$, $X = \sum_{i=0}^{\infty} x_i$ (well-defined since each coordinate becomes stationary) lies in $A$. 
Then, let $a_0 = X \in A = A_0$ and define $a_{n+1} = a_0 - \sum_{i=0}^n x_i$ for $n \ge 0$. Then, for $j \le n+1$,
$$(a_{n+1})_j = (a_0)_j - \left(\sum_{i=0}^n x_i\right)_j = \left(\sum_{i=0}^{\infty} x_i\right)_j - \left(\sum_{i=0}^n x_i\right)_j 
= \left(\sum_{i=0}^n x_i\right)_j - \left(\sum_{i=0}^n x_i\right)_j = 0,$$
so $a_{n+1} \in A_{n+1}$. We used the fact that for all $i \ge n+1 \ge j$, $(x_i)_j = 0$, thus reducing the sum to the first $n+1$ terms.
Finally, for any $n$, 
$$a_n - a_{n+1} = a_0 - \sum_{i=0}^{n-1} x_i - \left(a_0 - \sum_{i=0}^n x_i\right) = x_n,$$
so $(x_n) = d((a_i))$ and thus $(x_n) \in \im(d)$. Hence, $(x_n) \in \im(d)$ if and only if $\sum_{i=0}^{\infty}x_i \in A$.

Now, let $S : M \to \prod_{i=1}^{\infty} \Z/p\Z$ be the summation map defined above. Clearly, it is a $\Z$-module homomorphism. Then, we've shown
that for $z \in M$, $z \in \im(d)$ if and only if $S(z) \in \bigoplus_{i=1}^{\infty} \Z/p\Z$. Note that $S$ is surjective since for $y \in \prod_{i=1}^{\infty} \Z/p\Z$,
we can define $z \in M$ to be such that for all $n \ge 0$, $(z_n)_j = 0$ when $j \neq n+1$ and $(z_n)_{n+1} = y_{n+1}$. Then clearly, $S(z) = y$. Thus, the map
$$\overline{S} : M \to \frac{\prod_{i=1}^{\infty} \Z/p\Z}{\bigoplus_{i=1}^{\infty} \Z/p\Z}$$
is surjective with kernel equal to $\bigoplus_{i=1}^{\infty} \Z/p\Z$. We showed above that this is equal to $\im(d)$. Thus,
$\varprojlim^1 A_n = M/\im(d) \cong (\prod_{i=1}^{\infty} \Z/p\Z)(\bigoplus_{i=1}^{\infty} \Z/p\Z)$,
completing the proof.
\end{proof}

\begin{remark}
In retrospect, one could have simply noted that $d_A : \prod_{i=0}^{\infty} A \to \prod_{i=0}^{\infty} A$ is surjective so that we have an exact sequence
$A \to \varprojlim A/A_n \to \varprojlim^1 A_n \to 0$.
Then, from Exercise \ref{10.1}, $A = \bigoplus_{i=1}^{\infty}\Z/p\Z$ and $\varprojlim A/A_n \cong \prod_{i=1}^{\infty}\Z/p\Z$ giving us
$\varprojlim^1 A_n \cong (\prod_{i=1}^{\infty}\Z/p\Z)/(\bigoplus_{i=1}^{\infty}\Z/p\Z)$.
\end{remark}

\subsection{Exercise - Complete}\label{10.3}

\begin{proof}
Since $A$ is Noetherian and $M$ is a finitely generated $A$-module, we have by Krull's Theorem (Theorem 10.17) that $m \in \bigcap_{n=1}^{\infty} \mathfrak{a}^n M
= \ker(M \to \hat M)$ if and only if $m$ is annihilated by an element of $1 + \mathfrak{a}$.

Suppose $m \in \bigcap_{n=1}^{\infty} \mathfrak{a}^n M$ so that there is some $a \in \mathfrak{a}$
so that $(1 + a)m = 0$. Then, for any maximal ideal $\mathfrak{m}$ containing $\mathfrak{a}$, $1 + a \in 1 + \mathfrak{a} \subset 1 + \mathfrak{m}$. In particular, 
$1 + a \in A - \mathfrak{m}$ so that $1 + a$ is a unit in $A_{\mathfrak{m}}$. Since $(1+a)m = 0$ in $M$, we have the same relation in the $A_{\mathfrak{m}}$-module 
$M_{\mathfrak{m}}$. Because $1 + a$ is a unit in $A_{\mathfrak{m}}$, we have $0 = (1+a)^{-1}0 = (1+a)^{-1}(1+a)m = m$ so that $m \in \ker(M \to M_{\mathfrak{m}})$. Thus,
$m \in \bigcap_{\mathfrak{m} \supset \mathfrak{a}}\ker(M \to M_{\mathfrak{m}})$.

Conversely, let $N = \bigcap_{\mathfrak{m} \supset \mathfrak{a}}\ker(M \to M_{\mathfrak{m}})$. Then, $N_{\mathfrak{m}} = 0$ for all maximal ideals $\mathfrak{m}$
containing $\mathfrak{a}$. By Exercise \ref{3.14}, $N = \mathfrak{a}N$. Since $A$ is Noetherian, and $M$ is a finitely generated $A$-module, $M$ is Noetherian
by Proposition 6.5. Since $N$ is an $A$-submodule of $M$, $N$ is then finitely generated by Proposition 6.2. Thus, by Corollary 2.5, since $N = \mathfrak{a}N$ and
$N$ is finitely generated, there is some $x \in 1 + \mathfrak{a}$ such that $xN = 0$. Thus, every $m \in N$ is annihilated by $x \in 1 + \mathfrak{a}$, so 
$N \subset \bigcap_{n=1}^{\infty} \mathfrak{a}^n M$. Hence, we have equality:
$\bigcap_{\mathfrak{m} \supset \mathfrak{a}}\ker(M \to M_{\mathfrak{m}}) = \bigcap_{n=1}^{\infty} \mathfrak{a}^n M$.

If $\hat M = 0$, then $M = \ker(M \to \hat M) = \bigcap_{n=1}^{\infty} \mathfrak{a}^n M = \bigcap_{\mathfrak{m} \supset \mathfrak{a}}\ker(M \to M_{\mathfrak{m}})$. 
Thus, $M_{\mathfrak{m}} = 0$ for each $\mathfrak{m}$ containing $\mathfrak{a}$, so $\Supp(M)$ does not include
any $\mathfrak{m}$ over $\mathfrak{a}$. In fact, for any prime ideal $\mathfrak{p}$ such that $\mathfrak{m} \supset \mathfrak{p} \supset \mathfrak{a}$, the map
$M \to M_{\mathfrak{p}}$ factors as $M \to M_{\mathfrak{m}} \to M_{\mathfrak{p}}$, so $\ker(M \to M_{\mathfrak{p}}) \supset \ker(M \to M_{\mathfrak{m}}) = M$. Thus,
for any $\mathfrak{p} \in V(\mathfrak{a})$, $\mathfrak{p} \not\in \Supp(M)$, so $\Supp(M) \cap V(\mathfrak{a}) = \emptyset$.

Conversely, if $\Supp(M) \cap V(\mathfrak{a}) = \emptyset$, then $\ker(M \to M_{\mathfrak{m}}) = M$ for all maximal ideals $\mathfrak{m}$ containing $\mathfrak{a}$.
Thus, $M = \bigcap_{\mathfrak{m} \supset \mathfrak{a}}\ker(M \to M_{\mathfrak{m}}) = \bigcap_{n=1}^{\infty} \mathfrak{a}^n M = \ker(M \to \hat M)$. Thus, $\hat M = 0$.
Hence, $\hat M = 0$ if and only if $\Supp(M) \cap V(\mathfrak{a}) = \emptyset$.
\end{proof}

\subsection{Exercise - Complete}\label{10.4}

\begin{proof}
If $x \in A$ is not a zero divisor, then $A \xrightarrow{\cdot x} A$ is injective, thereby giving a short exact sequence
$$0 \to A \xrightarrow{\cdot x} A \to A/xA \to 0$$
of $A$-modules. Since $A$ is Noetherian, $A$ and thus $A/xA$ are Noetherian $A$-modules. Thus, each $A$-module in the sequence
is finitely generated (Proposition 6.2), so by Proposition 10.12, the sequence $0 \to \hat A \xrightarrow{\cdot \hat x} \hat A$
is exact. Hence, $\hat A \xrightarrow{\cdot \hat x} \hat A$ is injective, so $\hat x$ is not a zero divisor of $\hat A$.
\end{proof}

Even if $A$ is a Noetherian domain, $\hat A$ is not necessarily an integral domain.

\begin{proof}
Let $\mathfrak{a}$, $\mathfrak{b}$ be coprime ideals of a Noetherian domain $A$. Then, $\mathfrak{a}^n$ and $\mathfrak{b}^n$ are coprime for all $n \ge 1$. 
By Proposition 1.10 (Chinese Remainder Theorem), $A/(\mathfrak{ab})^n \cong A/\mathfrak{a}^n \times A/\mathfrak{b}^n$. Hence, we have a
commutative diagram
\begin{center}
\begin{tikzcd}
A/(\mathfrak{ab})^{n+1} \arrow[r, "\sim"] \arrow[d] & A/\mathfrak{a}^{n+1} \times A/\mathfrak{b}^{n+1} \arrow[d] \\
A/(\mathfrak{ab})^n \arrow[r, "\sim"] & A/\mathfrak{a}^n \times A/\mathfrak{b}^n.
\end{tikzcd}
\end{center}
For an ideal $I$, let $A^I$ denote the completion in the $I$-adic topology. Then, the inverse limit $A^{\mathfrak{a}\mathfrak{b}} = \varprojlim A/(\mathfrak{a}\mathfrak{b})^n$
is equal to $\varprojlim (A/\mathfrak{a}^n \times A/\mathfrak{b}^n)$ by Proposition 10.2. 

We have a map
$$\varprojlim (A/\mathfrak{a}^n \times A/\mathfrak{b}^n) \to \varprojlim A/\mathfrak{a}^n \times \varprojlim A/\mathfrak{b}^n$$
defined by the projections $A/\mathfrak{a}^n \times A/\mathfrak{b}^n \to A/\mathfrak{a}^n$ and $A/\mathfrak{a}^n \times A/\mathfrak{b}^n \to A/\mathfrak{b}^n$.
These are clearly compatible with the transition homomorphisms, so they induce the stated maps on inverse limits. Conversely, if $(x_n) \in \varprojlim A/\mathfrak{a}^n$
and $(y_n) \in \varprojlim A/\mathfrak{b}^n$, $(x_n,y_n)$ is a compatible sequence in $\prod_{n=1}^{\infty} (A/\mathfrak{a}^n \times A/\mathfrak{b}^n)$, giving 
the inverse map. Hence,
$$\varprojlim (A/\mathfrak{a}^n \times A/\mathfrak{b}^n) \xlongrightarrow{\sim} \varprojlim A/\mathfrak{a}^n \times \varprojlim A/\mathfrak{b}^n$$
so that $A^{\mathfrak{a}} \times A^{\mathfrak{b}} \cong \varprojlim A/\mathfrak{a}^n \times A/\mathfrak{b}^n \cong A^{\mathfrak{a}\mathfrak{b}}$.
Since $A^{\mathfrak{a}} \times A^{\mathfrak{b}}$ is not an integral domain, $A^{\mathfrak{a}\mathfrak{b}}$ is not an integral domain, thus
proving the counterexample.

For instance, $\Z$ is a Noetherian domain with coprime ideals $(2)$ and $(3)$, so with respect to the $(6)$-adic topology, we have $\hat \Z \cong \Z_2 \times \Z_3$,
which is not an integral domain, as $(1,0) \cdot (0,1) = 0$.
\end{proof}

\subsection{Exercise - Complete}\label{10.5}

\begin{proof}
For a Noetherian ring $A$, an ideal $I$, and an $A$-module $M$, let $M^I$ denote the $I$-adic completion. Suppose that $M$ is finitely
generated and $\mathfrak{a}$, $\mathfrak{b}$ are two ideals of $A$. Since $M$ is finitely generated and $A$ is Noetherian, $M$ is a 
Noetherian $A$-module (Proposition 6.5). Thus,
$$0 \to \mathfrak{b}^m M \to M \to M/\mathfrak{b}^m M \to 0$$
is a short exact sequence of finitely generated $A$-modules. Then, by Proposition 10.12, the sequence of $\mathfrak{a}$-adic completions
is also exact:
$$0 \to (\mathfrak{b}^m M)^{\mathfrak{a}} \to M^{\mathfrak{a}} \to (M/\mathfrak{b}^m M)^{\mathfrak{a}} \to 0.$$

By Proposition 10.13, if $N$ is any finitely generated $A$-module, then since $A$ is Noetherian, $N^{\mathfrak{a}} =  A^{\mathfrak{a}} \otimes_A N$.
Thus, the map $(\mathfrak{b}^m M)^{\mathfrak{a}} \to M^{\mathfrak{a}}$ factors as 
$$(\mathfrak{b}^m M)^{\mathfrak{a}} \xrightarrow{\sim} A^{\mathfrak{a}} \otimes_A (\mathfrak{b}^m M) \to A^{\mathfrak{a}} \otimes_A M \xrightarrow{\sim} \ M^{\mathfrak{a}}$$
and is given by $(b_m) \otimes x \mapsto (b_m) \otimes (x) \mapsto (b_m x)$. Therefore, the image of the map is $\mathfrak{b}^m M^{\mathfrak{a}}$. Then,
applying this to our exact sequence gives $(M/\mathfrak{b}^m M)^{\mathfrak{a}} \cong M^{\mathfrak{a}}/\mathfrak{b}^m M^{\mathfrak{a}}$.

Therefore, we have
$$(M^{\mathfrak{a}})^{\mathfrak{b}} = \varprojlim_m (M^{\mathfrak{a}}/\mathfrak{b}^m M^{\mathfrak{a}}) \cong \varprojlim_m (M/\mathfrak{b}^m M)^{\mathfrak{a}} 
\cong \varprojlim_m \varprojlim_n (M/\mathfrak{b}^m M)/\left(\frac{\mathfrak{a}^n M + \mathfrak{b}^m M}{\mathfrak{b}^m M}\right)$$
since $\mathfrak{a}^n(M/\mathfrak{b}^m M) \cong (\mathfrak{a}^n M + \mathfrak{b}^m M)/\mathfrak{b}^m M$ (Proposition 2.1) is the image of $\mathfrak{a}^n$ under the map 
$M \to M/\mathfrak{b}^m M$. By Proposition 2.1, 
$$(M/\mathfrak{b}^m M)/\left(\frac{\mathfrak{a}^n M + \mathfrak{b}^m M}{\mathfrak{b}^m M}\right) \cong M/(\mathfrak{a}^n M + \mathfrak{b}^m M),$$
which we will denote by $M_{n,m}$.

Now, note that an element $\varprojlim_n M_{n,m}$ gives coherent sequences of the form $x_{n+1,m} \mapsto x_{n,m}$. Then, $\varprojlim_m \varprojlim_n M_{n,m}$
gives coherent sequences $(x_m)$, where $(x_{n,m+1})_n = x_{m+1} \mapsto x_m = (x_{n,m})_n$ yields compatible, coherent maps $x_{n,m+1} \mapsto x_{n,m}$ for all
$n$. Then, if $n \le n'$ and $m \le m'$, we have a commutative diagram
\begin{center}
\begin{tikzcd}
M_{n',m'} \arrow[r]\arrow[d] & M_{n,m'} \arrow[d] \\
M_{n',m} \arrow[r] & M_{n,m}.
\end{tikzcd}
\end{center}
In particular, an element of the double inverse limit of $M_{n,m}$ is a doubly coherent sequence (with commutativity) $(x_{n,m})_{n,m}$. Then, given such a double inverse
limit, we can look at the diagonal $(x_{r,r})_r$, which is itself a coherent sequence. This defines a map $\varprojlim_m \varprojlim_n M_{n,m} \to \varprojlim_r M_{r,r}$.
Suppose $(x_{r,r})_r = 0$ so that $x_{r,r} = 0$ for all $r$. For all fixed $(n,m)$, there is some $r$ such that $n,m \le r$, so we have $x_{r,r} \mapsto x_{r,m} \mapsto x_{n,m}$
and thus $x_{n,m} = 0$. Hence, the map is injective. For surjectivity, suppose we have a coherent sequence $(y_r)_r \in \varprojlim_r M_{r,r}$. Define $x_{n,m}$ by
choosing some $r \ge n,m$; then, $x_{n,m}$ is the image of $y_r$ under $M_{r,r} \to M_{r,m} \to M_{n,m}$. By commutativity, this does not depend on $r$, so this defines
a coherent sequence in $\varprojlim_m \varprojlim_n M_{n,m}$, so the map is surjective and hence bijective. The map clearly preserves the $A$-module structure, so it is 
an isomorphism.

Therefore, we have
$$(M^{\mathfrak{a}})^{\mathfrak{b}} \cong \varprojlim_m \varprojlim_n M/(\mathfrak{a}^n M + \mathfrak{b}^m M) 
\cong \varprojlim_n M/(\mathfrak{a}^n M + \mathfrak{b}^n M) = \varprojlim_n M/(\mathfrak{a}^n + \mathfrak{b}^n) M.$$
Now, note that $(\mathfrak{a} + \mathfrak{b})^{2n} \subset \mathfrak{a}^n + \mathfrak{b}^n \subset (\mathfrak{a} + \mathfrak{b})^n$. Thus, the filtration 
$(\mathfrak{a}^n + \mathfrak{b}^n) M$ defines the same topology as $(\mathfrak{a} + \mathfrak{b})^n M$, so the completions are isomorphic. Hence,
$$(M^{\mathfrak{a}})^{\mathfrak{b}} \cong \varprojlim_n M/(\mathfrak{a}^n + \mathfrak{b}^n) M 
\cong \varprojlim_n M/(\mathfrak{a} + \mathfrak{b})^n M \cong M^{\mathfrak{a} + \mathfrak{b}},$$
as desired.
\end{proof}

\subsection{Exercise - Complete}\label{10.6}

\begin{proof}
Note that $\mathfrak{a}$ lies in the Jacobson radical if and only if $\mathfrak{m} \supset \mathfrak{a}$ for all maximal ideals
$\mathfrak{m}$ of $A$. Thus, it is sufficient to show that $\mathfrak{m}$ is closed in the $\mathfrak{a}$-topology if and only if
$\mathfrak{m} \supset \mathfrak{a}$.

If $\mathfrak{m} \supset \mathfrak{a}$, then for every $x \not\in \mathfrak{m}$, $x + \mathfrak{a} \subset x + \mathfrak{m}$ and
$(x + \mathfrak{m}) \cap \mathfrak{m} = \emptyset$. Since $x + \mathfrak{a}$ is open, their union over $x \not\in \mathfrak{m}$
is open. Note that
$$A - \mathfrak{m} \subset \bigcup_{x \not\in \mathfrak{m}} (x + \mathfrak{a}) 
\subset \bigcup_{x \not\in \mathfrak{m}} (x + \mathfrak{m}) \subset A - \mathfrak{m},$$
so $A - \mathfrak{m} = \bigcup_{x \not\in \mathfrak{m}} (x + \mathfrak{m})$ is open, and thus, $\mathfrak{m}$ is closed.

Suppose that $\mathfrak{m} \not\supset \mathfrak{a}$, so there is some $x \in \mathfrak{a} - \mathfrak{m}$. Since
$x \not\in \mathfrak{m}$, $x$ descends to a unit in $A/\mathfrak{m}$, so there is some $y \in A$ such that $\overline{x}\overline{y} = 1$.
Since $x \in \mathfrak{a}$, $xy \in \mathfrak{a}$ and $xy \in 1 + \mathfrak{m}$. Thus, replace $x$ with $xy$, so we have
$x \in \mathfrak{a} - \mathfrak{m}$ such that $x \in 1 + \mathfrak{m}$. Then, $x^n \in \mathfrak{a}^n$ and $x^n \in 1 + \mathfrak{m}$.
Thus, $1 - x^n \in (1 + \mathfrak{a}^n) \cap \mathfrak{m}$ for all $n > 0$. In the $\mathfrak{a}$-topology, $(1 - x^n)_{n>0}$ converges
to $1$. Since $1 \not\in \mathfrak{m}$, even though $(1 - x^n)_{n>0} \subset \mathfrak{m}$, $\mathfrak{m}$ is not closed.
\end{proof}

\begin{remark}
A Noetherian ring whose topology is defined by an ideal contained in the Jacobson radical is called a ``Zariski ring.''
\end{remark}

\subsection{Exercise - Complete}\label{10.7}

\begin{proof}
Since $A$ is a Noetherian ring, $\hat A$ is a flat $A$-algebra by Proposition 10.14. By Exercise \ref{3.16},
$A \to \hat A$ is faithfully flat if and only if for every $A$-module $M$, $M \to \hat A \otimes_A M$ is injective.
If $M \to \hat A \otimes_A M$ is not injective for some $A$-module $M$, then there is some nonzero $m \in M$ such that 
$m \mapsto 0$. Hence, $Am \to \hat A \otimes_A Am \to \hat A \otimes_A M$ is not injective. Since $\hat A$ is flat, and 
$Am \xhookrightarrow{} M$ is injective, $\hat A \otimes_A Am \to \hat A \otimes_A M$ is injective, so $Am \to \hat A \otimes_A Am$
is not injective. Thus, if there is an $A$-module $M$ such that $M \to \hat A \otimes_A M$ is not injective, there is a finitely
generated module $N$ such that $N \to \hat A \otimes_A N$ is not injective.

Therefore, to show that $M \to \hat A \otimes_A M$ is injective for every $A$-module $M$, it is sufficient to prove that 
$M \to \hat A \otimes_A M$ is injective for every 
finitely generated $A$-module $M$. In the case of finitely generated $A$-modules $M$, we have $\hat A \otimes_A M \cong \hat M$ 
since $A$ is Noetherian (Proposition 10.13). Therefore, we've reduced the proof to showing that $A$ is a Zariski ring if 
and only if $M \to \hat M$ is injective for all finitely generated $A$-modules $M$.

Suppose $A$ is a Zariski ring, i.e. the ideal $\mathfrak{a}$ defining its topology is contained in the Jacobson radical. Then,
for any finitely generated $A$-module $M$, we have by Corollary 10.19 that $\bigcap \mathfrak{a}^n M = 0$. This is the kernel 
of $M \to \hat M$ (Theorem 10.17), so $M \to \hat M$ is injective. 

Conversely, suppose for every finitely generated $A$-module
$M$, $M \to \hat M$ is injective. For each maximal ideal $\mathfrak{m}$ of $A$, $A/\mathfrak{m}$ is finitely generated, so 
$A/\mathfrak{m} \to \widehat{A/\mathfrak{m}}$ is injective. By Theorem 10.17, 
$\bigcap_n \mathfrak{a}^n (A/\mathfrak{m}) = 0$. Suppose to the contrary that $\mathfrak{m} \not\supset \mathfrak{a}$. Then,
there is some $x \in \mathfrak{a} - \mathfrak{m}$. Then, since $\mathfrak{m}$ is maximal, $A/\mathfrak{m}$ is a field. Since
$\overline{x} \neq 0$, there is some $y \in A$ such that $\overline{x} \overline{y} = 1$. Since $x \in \mathfrak{a}$, 
$z = xy \in \mathfrak{a}$ and $z \in 1 + \mathfrak{m}$. Thus, $z^n \in \mathfrak{a}$ and $z^n \equiv 1^n \equiv 1 \text{ mod }\mathfrak{m}$.
In particular, for every $n > 0$, $1 \in \mathfrak{a}^n(A/\mathfrak{m})$ since $z^n \cdot 1 = 1$, so 
$1 \in \bigcap_n \mathfrak{a}^n (A/\mathfrak{m}) = 0$, a contradiction. Thus, for all maximal ideals 
$\mathfrak{m}$ of $A$, $\mathfrak{a} \subset \mathfrak{m}$, so $\mathfrak{a}$ is contained in the Jacobson radical. 
Hence, $A$ is a Zariski ring.
\end{proof}

\subsection{Exercise}\label{10.8}

\subsection{Exercise}\label{10.9}

\subsection{Exercise}\label{10.10}

\subsection{Exercise}\label{10.11}

\subsection{Exercise}\label{10.12}
